\section{CSS Custom Properties (Variabel)}

\textbf{Custom properties} (variabel CSS) memungkinkan penyimpanan nilai yang dapat dipakai ulang di seluruh stylesheet. Dideklarasikan dengan sintaks \code{--nama: nilai;} dan diakses dengan fungsi \code{var(--nama)}. Biasanya didefinisikan pada \code{:root} agar tersedia secara global, tetapi dapat juga di-scope ke selector tertentu dan di-override di elemen turunan---mengikuti aturan cascade seperti properti CSS biasa \cite{mdnCss}.

\begin{webcode}[language=CSS, caption={Custom properties untuk design system}]
:root {
  --color-primary: #2a6fb0;
  --color-secondary: #1a4f80;
  --color-bg: #ffffff;
  --color-text: #333333;
  --spacing-sm: 0.5rem;
  --spacing-md: 1rem;
  --spacing-lg: 2rem;
  --radius: 8px;
}
/* Dark theme override */
[data-theme="dark"] {
  --color-bg: #1a1a2e;
  --color-text: #e0e0e0;
  --color-primary: #4da6ff;
}
.card {
  background: var(--color-bg);
  color: var(--color-text);
  padding: var(--spacing-md);
  border-radius: var(--radius);
  border: 1px solid var(--color-primary);
}
\end{webcode}

Custom properties memiliki beberapa keunggulan dibanding variabel preprocessor (Sass/Less): (1) bersifat \textbf{runtime}---dapat diubah melalui JavaScript atau media query; (2) mengikuti \textbf{cascade dan inheritance}---nilai berbeda di konteks berbeda; (3) mendukung \textbf{fallback}: \code{var(--warna, blue)} menggunakan \code{blue} jika variabel tidak terdefinisi. Ini menjadikannya ideal untuk tema dinamis (dark mode) yang dapat diubah tanpa reload halaman \cite{webdevCss}.

\begin{konsep}
\textbf{Design Token dengan Custom Properties.} Praktik modern mendefinisikan semua nilai desain (warna, spacing, tipografi, shadow, radius) sebagai custom properties di \code{:root}. Ini menciptakan \textit{design system} yang konsisten: mengubah satu variabel mengubah seluruh situs. Framework seperti Open Props menyediakan set variabel siap pakai. Untuk dark mode, cukup override variabel di \code{[data-theme="dark"]} atau \code{@media (prefers-color-scheme: dark)} \cite{cssTricks}.
\end{konsep}

